% Generated draft: requires mathematical and source review.
\documentclass[11pt]{article}
\usepackage[margin=25mm]{geometry}
\usepackage{fontspec}
\setmainfont{Latin Modern Roman}
\usepackage{amsmath,amssymb,mathtools,unicode-math}
\setmathfont{Latin Modern Math}
\usepackage{enumitem,needspace,xurl,hyperref}
\setlength{\parindent}{0pt}
\setlength{\parskip}{5pt}
\setlength{\emergencystretch}{3em}
\newcommand{\N}{\mathbb{N}}
\newcommand{\Z}{\mathbb{Z}}
\newcommand{\Q}{\mathbb{Q}}
\newcommand{\R}{\mathbb{R}}
\newcommand{\C}{\mathbb{C}}
\newcommand{\F}{\mathbb{F}}
\newcommand{\A}{\mathbb{A}}
\newcommand{\PP}{\mathbb{P}}
\newcommand{\E}{\mathbb{E}}
\newcommand{\eps}{\varepsilon}
\newcommand{\dd}{\,\mathrm{d}}
\begin{document}
\section*{UnsolvedMath ID: 30001314}
\textbf{Source problem code:} catalogue entry
\subsection{Definitions and mathematical statement}
For each integer \(g\geq 1\), let \(\mathcal H_g\) denote the source-intended class of properly embedded minimal helicoids of genus \(g\). The supplied record does not define this class sufficiently to determine its membership: in particular, it does not specify the ambient space or the end and asymptotic conditions meant by ``helicoid.'' The formulation below therefore makes the existence and uniqueness quantifiers precise, conditional on clarification of \(\mathcal H_g\).

In a fixed Euclidean ambient space \(\R^n\), an embedding is proper if the inverse image of every compact set is compact, and a smooth surface is minimal if its mean curvature vector vanishes identically. For surfaces \(M,N\subset\R^n\), define equivalence up to congruence and scaling by
\[
M\sim N\quad\Longleftrightarrow\quad
N=a+\lambda QM
\quad\text{for some }a\in\R^n,\ \lambda>0,\ Q\in O(n),
\]
where \(O(n)\) is the orthogonal group. The ambient dimension \(n\) here is a placeholder requiring source clarification, not an additional claim about the intended problem.

The normalized conjecture asks whether
\[
\forall g\in\Z_{\geq1}\quad
\exists M\in\mathcal H_g\quad
\forall N\in\mathcal H_g\quad N\sim M.
\]
Thus the assertion includes both existence for every positive integer genus and uniqueness up to the stated equivalence. The original statement attributes to Meeks the assertion that there is precisely one genus-\(g\) helicoid for every \(g\); the positive-genus range and the congruence-and-scaling qualification are explicit in the supplied normalized statement.

The original statement also contains a separate conjecture: there is a one-parameter family of Riemann-like properly embedded minimal surfaces with one handle inserted between two consecutive planar ends. This assertion is retained as source context, but the record supplies no definition of ``Riemann-like,'' parameter domain, notion of dependence on the parameter, or precise meaning of handle insertion. It therefore cannot be made into a fully specified additional conjecture from the supplied evidence alone.
\subsection{Short English statement}
The main question asks whether each positive integer genus admits exactly one properly embedded minimal helicoid, up to congruence and scaling. The intended class of helicoids needs clarification before this becomes a complete mathematical statement. The original source also includes a separate conjecture about a one-parameter family of Riemann-like surfaces with an inserted handle. Supplied dated curation reports partial progress in genus one; this draft does not independently verify that assessment.
\subsection{Sources}
\begin{itemize}
\item Local UnsolvedMath registry: \texttt{data/unsolvedmath/problems.json}, permanent ID \(30001314\), title ``Uniqueness of Genus-\(g\) Helicoids,'' exact source problem code \texttt{catalogue entry}. This draft uses the supplied complete registry record.
\item Supplied report citation: \emph{Differentialgeometrie im Großen (2010)}, Oberwolfach Reports, pp.\ 1943--1972, DOI \texttt{10.4171/owr/2009/35}. Supplied literature-source label: ``Differential geometry report.'' \url{https://doi.org/10.4171/owr/2009/35}.
\item Attribution: the original statement attributes the genus-\(g\) helicoid conjecture to Meeks. The separate Riemann-like-family conjecture has no proposer identified in the supplied record. The registry gives proposed year 2009.
\item Supplied literature source: ``Embedded genus-one helicoid.'' \url{https://doi.org/10.4007/annals.2009.169.347}.
\item Supplied literature source: ``Bernstein--Breiner symmetry theorem.'' \url{https://arxiv.org/abs/1011.3769}.
\item Literature assessment dated exactly \texttt{2026-08-21}: the supplied curation labels the problem partially solved and reports existence and embeddedness of a genus-one helicoid, proved symmetry and asymptotic structure, and an open genus-one uniqueness assertion. This is supplied dated curation, not an independently verified literature assessment; the record itself requests independent verification of time-sensitive status and cited claims.
\item The registry labels the normalized statement \texttt{corrected\_verified} and reports that its mathematical content was checked against the official Oberwolfach report, with notation, context, and question form normalized. This verification claim is supplied provenance and has not been independently checked for this draft.
\end{itemize}
\end{document}
